\documentclass[a4paper,12pt,oneside]{article}

\usepackage[T2A]{fontenc}
\usepackage[utf8]{inputenc}
\usepackage[russian]{babel}

\usepackage{lmodern}
\usepackage{microtype}

\usepackage{amsmath,amssymb,amsthm,mathtools}

\usepackage{geometry}
\geometry{top=25mm,left=25mm,right=20mm,bottom=25mm}

\usepackage{enumitem}
\setlist[enumerate]{leftmargin=*, itemsep=0.35\baselineskip, topsep=0.35\baselineskip}

\newlist{alphlist}{enumerate}{1}
\setlist[alphlist]{label=\textbf{\alph*)}, leftmargin=*, align=left}

\newlist{numlist}{enumerate}{1}
\setlist[numlist]{label=\textbf{\arabic*)}, leftmargin=*, align=left}

\newcommand{\Task}[1]{%
  \subsection*{Задача #1}%
  \vspace{-0.2\baselineskip}%
}

\setlength{\parindent}{1.25em}
\setlength{\parskip}{0pt}

\begin{document}

\begin{center}
    {\Large\bfseries Бинарные отношения.}\\[2mm]
\end{center}
\vspace{2mm}
\hrule
\vspace{4mm}

\Task{1. Найдите область определения и множество значений отношений}
\begin{alphlist}
  \item $\bigl\{(a,1),(a,2),(a,3),(a,4),(a,5),(a,6)\bigr\}$
  \item $\bigl\{(x,y)\,:\,x,y\in\mathbb{Z},\;x^2+y^2\le 16\bigr\}$
  \item $\bigl\{(x,y)\,:\,0\le x\le 10,\;0\le y\le 10,\;x>2y\bigr\}$
\end{alphlist}

\Task{2. Пусть отношения $U, V \subseteq \mathbb{R}\times\mathbb{R}$ определены указанным ниже способом}
\[
U=\bigl\{(x,y)\,:\,y=x^2+5\bigr\},
\qquad
V=\bigl\{(x,y)\,:\,y=3x\bigr\}.
\]
\begin{alphlist}
  \item Опишите отношение $V\circ U$.
  \item Опишите отношение $U\circ V$.
  \item Опишите отношение $U^{-1}$.
  \item Опишите отношение $V^{-1}$.
  \item Найдите множество значений (образ) отношения $U$.
\end{alphlist}

\Task{3}
$A=\{a,b,c,d,e\}$. Опишите наименьшее рефлексивное отношение на этом множестве.

\Task{4}
\begin{numlist}
  \item Множество вершин $\{a,b,c,d,e,f\}$ и отношение $R$ для рёбер имеет вид:
  \[
    R_1=\bigl\{(a,b),(b,c),(a,c),(b,e),(c,f),(c,d),(d,f),(f,e)\bigr\}.
  \]

  \item Множество вершин $\{a,b,c,d\}$ и отношение $R$ для рёбер имеет вид:
  \[
    R_2=\bigl\{(b,c),(a,d),(b,a),(d,c),(b,d),(c,a),(d,d)\bigr\}.
  \]

  \item Определите свойства $R_1$ и $R_2$.
\end{numlist}

\Task{5}
Пусть $A=\{1,2,3\}$, а отношение $R$ на $A$ задано упорядоченными парами:
\[
R=\bigl\{(1,1),(2,1),(1,3),(3,1),(2,3)\bigr\}.
\]
Найдите рефлексивное, симметричное и транзитивное замыкания.

\Task{6}
Отношение $R_3$ задано на множестве $A=\{1,2,3,4,5\}$:
\[
a\,R_3\,b \iff a+b \text{ чётное}.
\]
Чему равна мощность множества классов эквивалентности $R_3$?

\Task{7}
Отношение $R$ на вещественной прямой $\mathbb{R}$ задано условием:
\[
xRy \iff x-y \text{ — целое число}.
\]
Докажите, что $R$ — отношение эквивалентности, и опишите классы эквивалентности,
содержащие $0$, $\tfrac{1}{2}$ и $\sqrt{2}$.

\end{document}
